\begin{comment}	
O câncer de próstata é uma das neoplasias malignas mais comuns em homens, representando um desafio significativo para a saúde pública \cite{ref_siegel2024}. Estatísticas coletadas pela \textit{International Agency for Research on Cancer} indicam que, em 2022, essa neoplasia foi a quarta mais comum no mundo, correspondendo a 7,3\% de todos os casos de câncer \cite{JournalStatistic2024}. Neste contexto, a gradação tumoral desempenha papel central na definição do manejo clínico e do prognóstico. O sistema de \textit{Gleason}, baseado na soma dos dois padrões histológicos predominantes, evoluiu para o sistema de graduação proposto pela \textit{International Society of Urological Pathology} (ISUP), organizado em cinco grupos que variam de 1 a 5. Este sistema fornece uma estratificação prognóstica aprimorada e é amplamente utilizado na prática clínica \cite{ref-epstein2016}.


A avaliação histopatológica tradicional de imagens de lâminas inteiras (\textit{Whole-Slide Images} - WSIs) depende da análise manual realizada por patologistas. Neste sentido, tendo em vista a quantidade de campos a serem observados em uma única lâmina, o processo de análise do exame torna-se uma tarefa repetitiva e exige jornadas de trabalho prolongadas, foco contínuo e especialização do profissional. Nesse sentido, a fadiga decorrente da análise contínua durante esse processo pode introduzir variabilidades inter e intraobservador nos resultados, o que influencia diretamente o manejo clínico do paciente~\cite{ref-bulten2020}. Com isso, nos últimos anos, tem sido observado um crescimento significativo no uso de técnicas de Inteligência Artificial (IA) para automatizar a análise de exames histopatológicos e auxiliar no diagnóstico do câncer de próstata, com o intuito de mitigar os problemas citados anteriormente \cite{expert2022review}.

Neste contexto, as Redes Neurais Convolucionais (RNCs) têm representado o \textit{estado-da-arte} na análise de imagens médicas \cite{ref-campanella2019}, sendo capazes de extrair características relevantes e combinar padrões visuais complexos para identificar determinadas condições patológicas \cite{araujo2021active}. Tal capacidade torna essa abordagem particularmente promissora para a análise automatizada dos padrões de Gleason e dos grupos de grau ISUP. Com isso, a abordagem contribui para a padronização do diagnóstico do câncer de próstata e para a redução das variabilidades intra e interobservador, inerentes à avaliação histopatológica convencional. Assim, além de reduzir o tempo de diagnóstico, a abordagem viabilizaria um manejo clínico mais preciso e reprodutível.

Este trabalho propõe um pipeline estruturado para a classificação automática dos grupos de grau ISUP em imagens histopatológicas de próstata. As principais contribuições são: (i) uma função de perda híbrida que combina penalização ordinal e perda focal, considerando a natureza hierárquica do problema e o desbalanceamento da base de dados; (ii) um método automatizado de remoção de ruído baseado na incerteza estimada pela média das predições dos modelos EfficientNet (B0, B1, B2 e B3); e (iii) a avaliação comparativa das arquiteturas EfficientNet B0, B3 e B7.	conteúdo...
\end{comment}

